\section{Related Work}
\label{sec:rel_work}

%==========================================================================
% PÁRRAFO 1: Limitaciones de conexiones horizontales en V1
%==========================================================================
% Early computational approaches to extra-classical receptive field (eCRF) effects primarily emphasized the role of intrinsic horizontal connections within the primary visual cortex (V1). These lateral projections allow for spatial integration by linking neurons with non-overlapping receptive fields across the cortical surface. However, horizontal connections face significant physiological constraints regarding their spatial extent and conduction velocity, which often fail to match the rapid, long-range dynamics of surround suppression observed experimentally. While lateral inhibition within V1 can produce a measurable decrease in firing rates, simulations often show that it typically yields a modest modulation---frequently failing to exceed a 30--35\% reduction in activity \commentcd{reference}. This is insufficient to reach the robust levels of suppression ($\sim$75\%) reported in biological systems, suggesting that local V1 circuitry alone lacks the necessary dynamic range to fully account for global contextual integration.

%==========================================================================
% PÁRRAFO 2: Limitaciones del feedback lineal simple
%==========================================================================
% \commentcd{FUSIONAR ESTE PARRAFO CON EL ANTERIOR} While hierarchical feedback from higher-order areas such as V2 has been widely proposed as a driver of long-range contextual modulation, conventional computational implementations often rely on simplified feedback loops. These models typically treat top-down signals as linear additive inputs or non-specific gain modulators, which fail to capture the highly non-linear and stimulus-specific nature of extra-classical receptive field (eCRF) effects. Crucially, without internal competitive dynamics within the feedback-generating areas, top-down signals cannot achieve the precise spatial and feature-specific filtering required to match physiological observations. This limitation prevents existing architectures from reaching robust levels of surround suppression ($\sim$75\%), underscoring the necessity of a framework where hierarchical competition refines the feedback signals before they influence the laminar microcircuits of V1.

Early computational approaches to eCRF effects primarily emphasized the role of intrinsic horizontal connections within the primary visual cortex (V1) \commentcd{reference}. These lateral projections allow for spatial integration by linking neurons with non-overlapping receptive fields across the cortical surface. However, horizontal connections face significant physiological constraints regarding their spatial extent and conduction velocity, which often fail to match the rapid, long-range dynamics of surround suppression observed experimentally. Moreover, simulations of lateral inhibition within V1  yields modest modulation—around 30--35\% reduction in neuron firing rate \commentcd{reference}. Far from the ($\sim$75\%) suppression reported in biological systems. On the other hand, hierarchical feedback from higher-order areas such as V2 has also been proposed as a driver of long-range contextual modulation. Current computational implementations rely on simplified feedback loops between V1 and V2 \commentcd{eso que dije en la frase anterior es cierto}. These models typically treat top-down signals as linear additive inputs or non-specific gain modulators, which fail to capture the highly non-linear and stimulus-specific nature of eCRF effects. Crucially, without internal competitive dynamics within the feedback-generating areas, top-down signals cannot achieve the precise spatial and feature-specific filtering required to match physiological observations. \commentcd{me parece que esta frase adelanta los resultados, o es algo que ya se sabía y podemos poner una referencia?}

%==========================================================================
% TODO: PÁRRAFO NUEVO 3 - Modelos de normalización divisiva
%==========================================================================
% [AQUÍ FALTA: Conectar con la literatura de normalización divisiva, tanto
% biológica como artificial:
%   - Normalización divisiva como operación canónica (Carandini & Heeger, 2012)
%   - Modelos clásicos: Heeger (1992), modelo de contraste de Carandini et al.
%   - Implementaciones en deep learning: batch norm, layer norm, divisive
%     normalization networks
%   - Cómo estos modelos NO incorporan competencia jerárquica entre áreas
%

In the domain of spiking neural networks, considerable progress has been made in developing biologically constrained models of cortical processing. The canonical microcircuit model of \citet{Potjans2014} has emerged as a gold standard for reproducing the laminar connectivity and asynchronous-irregular dynamics of local cortical patches. This model has been successfully extended to multi-area hierarchies \citep{Schmidt2018a, Schmidt2018b} and applied to cognitive tasks such as working memory and decision-making. Concurrently, spiking networks trained via surrogate gradients have achieved competitive performance on visual recognition benchmarks \citep{Shrestha2018}, demonstrating that spike-based computation can scale to complex tasks. However, these models typically treat hierarchical interactions as feedforward feature extraction or simple convergent feedback, without incorporating lateral inhibitory competition within higher areas. \commentcd{ De nuevo me entra la duda si esta frase que sigue les estamos adelantando el resultado o si es algo que se sabe y podemos citar} Consequently, they do not exhibit the robust, stimulus-specific surround suppression that characterizes biological vision, nor do they leverage the computational advantages of hierarchical winner-take-all dynamics for context-dependent normalization.

%==========================================================================
% TODO: PÁRRAFO NUEVO 4 - SNNs bio-inspiradas y modelos jerárquicos
%==========================================================================
% [AQUÍ FALTA: Revisar modelos de SNN existentes para procesamiento visual
% y mostrar cómo ninguno implementa competencia lateral en áreas superiores:
%   - SNNs para tareas de visión (Diehl & Cook, 2015; Shrestha & Orchard, 2018)
%   - Modelos de microcircuito de Potjans & Diesmann (2014) como gold standard
%   - Extensiones jerárquicas existentes del microcircuito (si las hay)
%   - Qué falta en todos ellos: inhibición lateral en V2
%
The phenomenon of surround suppression has been extensively studied through the lens of divisive normalization, a canonical neural computation in which a neuron's response is divided by the aggregate activity of a broader neuronal pool \citep{Carandini2012}. Early formulations by \citet{Heeger1992} demonstrated that divisive normalization can account for contrast-dependent saturation and cross-orientation suppression in V1 using simple computational rules. Subsequent work generalized this principle to a wide range of sensory and cognitive phenomena \citep{Carandini2012, Reynolds2009}. In parallel, the machine learning community has independently converged on normalization as an essential architectural component, with batch normalization \citep{Ioffe2015}, layer normalization \citep{Ba2016}, and divisive normalization networks \citep{Ren2017} all demonstrating that rescaling activations relative to local or global statistics improves training stability and generalization. However, these implementations typically operate through feedforward or self-normalizing mechanisms within a single processing stage, lacking the inter-areal competitive dynamics that characterize cortical hierarchies.

%==========================================================================
% TODO: PÁRRAFO NUEVO 5 - Predictive coding y hierarchical inference
%==========================================================================
% [AQUÍ FALTA: Conectar mecanismo de WTA en V2 con predictive coding y
% explaining away:
%   - Rao & Ballard (1999): predictive coding en corteza visual
%   - Friston (2005): free-energy principle y hierarchical inference
%   - Explicar cómo tu WTA en V2 es un mecanismo biológicamente plausible de
%     "explaining away"
%   - Relación con modelos de atención y sparse coding
%
On this work, we propose a hierarchical competition mechanism to explain the eCRF effects. This mechanisms is closely aligned with predictive coding frameworks \citep{Rao1999, Friston2005}, where higher cortical areas generate top-down predictions that are compared against incoming sensory evidence at \commentcd{sera esto un from en vez de at} lower levels; mismatches produce prediction errors that propagate upward, while accurately predicted features are suppressed through a process known as \emph{explaining away}. This computation bears functional similarities to biased competition models of attention \citep{Desimone1995, Reynolds2009}, where top-down signals resolve competition among multiple stimuli by selectively enhancing or suppressing specific feature representations. The WTA competition that we implemented in V2 can be interpreted as a circuit-level instantiation of these principle: when a peripheral stimulus drives a competing V2 module, it actively inhibits the center-aligned module's feedback, effectively signaling that the central activity is redundant given the broader spatial context.  By grounding these abstract theories in a concrete spiking microcircuit with laminar specificity, our work provides a bridge between normative computational principles and their biologically plausible implementation.
%==========================================================================
% TODO: PÁRRAFO NUEVO 6 - Resumen del gap y cómo lo abordamos
%==========================================================================
% [AQUÍ FALTA: Un cierre que sintetice el gap en la literatura y posicione
% el trabajo como la solución. No más de 3-4 oraciones.
%

In summary, non of  these research lines --- biophysical models of cortical microcircuits, divisive normalization in theory and practice, and predictive coding ---  explains how extra-classical surround suppression emerges from hierarchical interactions. By grounding these abstract theories in a concrete spiking microcircuit with laminar specificity, our work provides a bridge between normative computational principles and their biologically plausible implementation. Specifically, the role of lateral inhibitory competition in higher visual areas as a gating mechanism for top-down feedback remains unexplored in large-scale spiking models. Our work addresses this gap and demonstrates that V2 lateral inhibition is both necessary and \commentcd{entiendo que es necesaria, pero no por qué es suficiente} sufficient to reproduce the hallmark signatures of robust sensory normalization.